\section*{Pregunta 1}
\noindent Realice un informe sobre:

\begin{itemize}
  \item[a] ¿Qué es un prototipo de un producto?

  \item[b] ¿Cuál es la utilidad de un prototipo?

  \item[c] ¿Qué objetivos se tiene al realizar un prototipo de un producto?
\end{itemize}

\subsection*{¿Qué es un prototipo de un producto?}
\noindent Un prototipo es una versión preliminar o modelo de un
producto, servicio o sistema. Se trata de la primera representación
tangible de una idea; algo que se puede ver, sostener y probar antes
de que el producto entre en producción.

\noindent El prototipo no es el producto final, sino una herramienta
de aprendizaje y validación. Puede variar en su grado de detalle y
funcionalidad, desde un simple boceto en papel hasta un modelo casi
idéntico al producto terminado. Su propósito principal es simular,
validar o probar la funcionalidad, el diseño o la viabilidad de una
idea antes de invertir recursos significativos en su desarrollo final.

\noindent Existen diferentes tipos de prototipos, que se clasifican
según su "fidelidad" o nivel de detalle:

\begin{itemize}
  \item Prototipos de baja fidelidad: Son representaciones rápidas y
    sencillas, como bocetos en papel, maquetas con materiales simples
    o storyboards. Se utilizan para probar conceptos básicos y
    obtener retroalimentación temprana.

  \item Prototipos de media fidelidad: Son esquemas más
    estructurados, como wireframes digitales o mockups, que definen
    la arquitectura y la lógica de interacción, pero aún carecen del
    diseño visual final.

  \item Prototipos de alta fidelidad: Se acercan mucho al producto
    final, incluyendo colores, tipografía, interactividad completa y,
    en muchos casos, funcionalidad real. Se utilizan para pruebas de
    usuario exhaustivas y presentaciones a inversores.
\end{itemize}

\subsection*{¿Cuál es la utilidad de un prototipo?}
\noindent La creación de un prototipo aporta un valor inmenso al
proceso de desarrollo de productos, ofreciendo múltiples utilidades prácticas:

\begin{itemize}
  \item Detección temprana de errores: Permite identificar fallos de
    diseño, problemas de funcionamiento o defectos en etapas
    iniciales, cuando son más fáciles y económicos de corregir. Esto
    evita modificaciones costosas después de que el producto ya ha
    sido fabricado en grandes cantidades.

  \item Validación con el usuario: Facilita la recopilación de
    comentarios reales del público objetivo. Al interactuar con el
    prototipo, los usuarios potenciales pueden expresar sus
    impresiones, lo que ayuda a refinar el producto para que
    satisfaga sus necesidades.

  \item Visualización y comunicación: Convierte una idea abstracta en
    algo tangible, lo que mejora la comunicación entre los equipos de
    diseño, ingeniería, producción y marketing. Un prototipo permite
    "mostrar" en lugar de "contar", facilitando la discusión y la
    toma de decisiones.

  \item Reducción de riesgos: Ayuda a minimizar los riesgos asociados
    con el desarrollo de un nuevo producto. Al probar la viabilidad
    técnica y la aceptación del mercado antes de la producción en
    masa, se reduce la probabilidad de invertir en un producto que no
    funcionará o no se venderá.

  \item Atracción de inversores: Un prototipo físico o funcional es
    una herramienta poderosa para presentar la idea a inversores
    potenciales, ya que demuestra el concepto de una manera mucho más
    convincente que un simple plan de negocios.
\end{itemize}

\subsection*{¿Qué objetivos se tiene al realizar un prototipo de un producto?}
\noindent Los objetivos de la creación de un prototipo se alinean
directamente con sus utilidades. Los principales son:

\begin{itemize}
  \item Validar la viabilidad y funcionalidad: Proporcionar
    evidencias sobre si una idea es técnica y funcionalmente
    factible. El objetivo es comprobar si la idea tiene sentido y
    responde a las necesidades detectadas.

  \item Probar y refinar el diseño: Detectar errores de diseño o
    usabilidad específicos para encontrar soluciones y desarrollar un
    producto que se ajuste a las necesidades de los usuarios.

  \item Obtener retroalimentación y aprender: El objetivo principal
    es aprender de forma rápida y con el menor coste posible. Cada
    iteración del prototipo ofrece una oportunidad para mejorar y
    acercarse a la solución óptima.

  \item Visualizar y comunicar la idea: Tangibilizar conceptos
    abstractos para que los usuarios, clientes o el equipo puedan
    imaginar el resultado final.

  \item Minimizar riesgos y maximizar la rentabilidad: Al identificar
    y solucionar problemas antes de la producción, se ahorra tiempo y
    dinero a largo plazo, lo que contribuye a maximizar la
    rentabilidad del producto.
\end{itemize}
